\chapter{Những Điều Về Cuộc Sống}
\markboth{Những Điều Về Cuộc Sống}{Things a Teacher Wants to Say About Life}

\section{Thời Gian Rất Quý}
\begin{truyen}{Thời Gian Rất Quý}{Time Is Precious}
\chuhoa{T}{hời gian đi rất đều, dù em có dùng nó tử tế hay không.}
\textit{(Time moves steadily whether you use it wisely or not.)}

Điều làm nó quý là một ngày qua rồi thì không lấy lại được.
\textit{(What makes it precious is that once a day passes, it cannot be reclaimed.)}

Học sinh thường nghĩ mình còn nhiều thời gian lắm.
\textit{(Students often think they have plenty of time.)}

Đúng là còn trẻ, nhưng chính vì vậy càng nên tập dùng thời gian cho đàng hoàng từ sớm.
\textit{(It is true that youth has time, but that is exactly why it should be used well from early on.)}
\end{truyen}

\begin{baihoc}
Tôn trọng thời gian là một cách tôn trọng chính cuộc đời mình.
\textit{(Respecting time is one way of respecting your own life.)}
\end{baihoc}

\begin{ghinhoanh}
\textbf{Short English to remember:}

Time is valuable because it does not return.

\end{ghinhoanh}

\begin{tuhocnhanh}
1. Vì sao thời gian quý? \textit{Why is time precious?}

2. Em đang lãng phí thời gian nhiều nhất ở đâu? \textit{Where are you wasting the most time?}

3. Em có thể dùng ngày mai đỡ phí hơn thế nào? \textit{How can you use tomorrow more wisely?}
\end{tuhocnhanh}
\ngancach

\section{Sức Khỏe Quan Trọng}
\begin{truyen}{Sức Khỏe Quan Trọng}{Health Matters}
\chuhoa{N}{hiều học sinh chỉ nhớ sức khỏe quan trọng khi cơ thể bắt đầu trả giá.}
\textit{(Many students remember health matters only when the body begins to pay the price.)}

Ngủ thiếu, ăn thất thường, căng kéo dài rồi vẫn cố coi đó là bình thường.
\textit{(Too little sleep, irregular eating, long stress, and still treating that as normal.)}

Nhưng đầu óc muốn học tốt cũng cần một cơ thể được đối xử đàng hoàng.
\textit{(But a mind that wants to learn well also needs a body treated properly.)}

Giữ sức khỏe không làm em chậm lại. Nó giữ cho em còn sức để đi đường dài.
\textit{(Protecting health does not slow you down. It helps you endure the long road.)}
\end{truyen}

\begin{baihoc}
Cơ thể không phải cái máy để em bắt nó gồng mãi mà không trả giá.
\textit{(The body is not a machine you can force forever without cost.)}
\end{baihoc}

\begin{ghinhoanh}
\textbf{Short English to remember:}

A healthy body supports a learning mind.

\end{ghinhoanh}

\begin{tuhocnhanh}
1. Vì sao sức khỏe ảnh hưởng việc học? \textit{Why does health affect learning?}

2. Cơ thể em đang báo điều gì mà em hay bỏ qua? \textit{What is your body telling you that you often ignore?}

3. Em cần sửa thói quen nào để khỏe hơn? \textit{What habit do you need to change for better health?}
\end{tuhocnhanh}
\ngancach

\section{Kiến Thức Là Tài Sản}
\begin{truyen}{Kiến Thức Là Tài Sản}{Knowledge Is an Asset}
\chuhoa{K}{iến thức là thứ đi cùng em lâu hơn nhiều món đồ khác.}
\textit{(Knowledge stays with you much longer than many other possessions.)}

Nó giúp em hiểu việc, hiểu người, hiểu đời và tự bảo vệ mình.
\textit{(It helps you understand work, people, life, and protect yourself.)}

Có tài sản mất vì trộm cắp hay biến động.
\textit{(Some assets can be lost through theft or instability.)}

Nhưng điều em đã học thật sự và giữ được trong đầu thường là một dạng của cải rất bền.
\textit{(But what you have truly learned and kept is a very durable form of wealth.)}
\end{truyen}

\begin{baihoc}
Đầu óc được bồi đắp tốt là một tài sản rất khó ai lấy khỏi em.
\textit{(A well-developed mind is an asset that is hard for anyone to take away.)}
\end{baihoc}

\begin{ghinhoanh}
\textbf{Short English to remember:}

Real knowledge is durable wealth.

\end{ghinhoanh}

\begin{tuhocnhanh}
1. Vì sao kiến thức được gọi là tài sản? \textit{Why can knowledge be called an asset?}

2. Nó giúp em bảo vệ mình bằng cách nào? \textit{How does knowledge help protect you?}

3. Em đang tích lũy tài sản đầu óc của mình ra sao? \textit{How are you building your mental wealth?}
\end{tuhocnhanh}
\ngancach

\section{Đọc Sách Mở Rộng Thế Giới}
\begin{truyen}{Đọc Sách Mở Rộng Thế Giới}{Books Widen the World}
\chuhoa{S}{ách làm em đi xa hơn nơi mình đang sống.}
\textit{(Books carry you beyond the place where you live.)}

Nó cho em gặp những con người khác, cách nghĩ khác, nỗi đau khác và cả những hy vọng khác.
\textit{(They let you meet different people, thoughts, pains, and hopes.)}

Một người đọc lâu năm thường khó sống quá hẹp.
\textit{(A long-term reader often finds it harder to live too narrowly.)}

Đọc sách là một trong những cách yên lặng nhưng mạnh để mở rộng thế giới bên trong mình.
\textit{(Reading is one of the quiet but strong ways to widen the inner world.)}
\end{truyen}

\begin{baihoc}
Đọc không chỉ cho thêm thông tin. Nó còn giúp con người bớt nông và bớt hẹp.
\textit{(Reading does not only add information. It also helps people become less shallow and less narrow.)}
\end{baihoc}

\begin{ghinhoanh}
\textbf{Short English to remember:}

Books widen both thought and heart.

\end{ghinhoanh}

\begin{tuhocnhanh}
1. Vì sao sách giúp mở rộng thế giới? \textit{Why do books widen the world?}

2. Người đọc lâu năm thường khác ở điểm nào? \textit{How are long-term readers often different?}

3. Cuốn sách nào từng làm em nhìn rộng hơn? \textit{What book has widened your view?}
\end{tuhocnhanh}
\ngancach

\section{Học Suốt Đời}
\begin{truyen}{Học Suốt Đời}{Keep Learning for Life}
\chuhoa{R}{a khỏi trường không có nghĩa là việc học kết thúc.}
\textit{(Leaving school does not mean learning ends.)}

Đời thay đổi, công việc đổi, con người cũng đổi.
\textit{(Life changes, work changes, and people change too.)}

Ai ngừng học quá sớm thường bị đời bỏ lại rất nhanh.
\textit{(Those who stop learning too early are often left behind by life quickly.)}

Học suốt đời không phải lúc nào cũng là ngồi lớp. Nó là giữ đầu óc còn mở, còn muốn hiểu thêm và sửa thêm.
\textit{(Lifelong learning is not always classroom study. It is keeping the mind open, still wanting to understand and improve.)}
\end{truyen}

\begin{baihoc}
Người còn học được thì còn lớn lên được.
\textit{(Those who can still learn can still grow.)}
\end{baihoc}

\begin{ghinhoanh}
\textbf{Short English to remember:}

Learning should outlive school.

\end{ghinhoanh}

\begin{tuhocnhanh}
1. Vì sao ra khỏi trường không có nghĩa là hết học? \textit{Why does leaving school not mean the end of learning?}

2. Học suốt đời thể hiện bằng những cách nào? \textit{How does lifelong learning show itself?}

3. Em có muốn giữ đầu óc học tiếp sau này không? \textit{Do you want to keep learning later in life?}
\end{tuhocnhanh}
\ngancach

\section{Cuộc Đời Không Dễ Dàng}
\begin{truyen}{Cuộc Đời Không Dễ Dàng}{Life Is Not Easy}
\chuhoa{K}{hông phải để làm em sợ, nhưng thầy phải nói thật: cuộc đời không dễ.}
\textit{(Not to frighten you, but I have to say honestly: life is not easy.)}

Sẽ có lúc em vất vả, hiểu lầm, thất bại, hoặc phải đi qua những đoạn rất không giống điều mình mong.
\textit{(There will be times of hardship, misunderstanding, failure, and roads unlike what you hoped for.)}

Nhưng biết đời không dễ không phải để bi quan. Nó để em bớt ảo tưởng và chuẩn bị tốt hơn.
\textit{(Knowing life is not easy is not for pessimism. It helps you abandon illusion and prepare better.)}

Người chuẩn bị tinh thần đúng thường đỡ gãy hơn khi va phải sự thật.
\textit{(Those who prepare their minds honestly break less when reality hits.)}
\end{truyen}

\begin{baihoc}
Sự thật rằng đời không dễ không làm mất hy vọng. Nó chỉ làm hy vọng của em bớt ngây thơ hơn.
\textit{(The truth that life is not easy does not kill hope. It makes hope less naive.)}
\end{baihoc}

\begin{ghinhoanh}
\textbf{Short English to remember:}

Honest preparation makes hardship less shocking.

\end{ghinhoanh}

\begin{tuhocnhanh}
1. Vì sao cần nhìn thật rằng đời không dễ? \textit{Why is it important to see honestly that life is not easy?}

2. Biết điều đó giúp mình chuẩn bị ra sao? \textit{How does that truth help us prepare?}

3. Em đang hy vọng quá ngây thơ ở điều gì? \textit{Where is your hope still too naive?}
\end{tuhocnhanh}
\ngancach

\section{Người Tốt Luôn Cần Thiết}
\begin{truyen}{Người Tốt Luôn Cần Thiết}{Good People Are Always Needed}
\chuhoa{T}{hế giới có thể cần người giỏi, nhưng nó cũng rất cần người tốt.}
\textit{(The world may need skilled people, but it also deeply needs good people.)}

Người tốt không phải người yếu.
\textit{(A good person is not a weak person.)}

Đó là người giữ được phần tử tế trong cách đối xử, trong cách làm việc, trong cách dùng khả năng của mình.
\textit{(It is someone who keeps decency in behavior, work, and the use of their abilities.)}

Nếu chỉ giỏi mà làm người khác đau thêm, sự giỏi ấy không đẹp trọn.
\textit{(If someone is skilled but harms others, that skill is not fully beautiful.)}
\end{truyen}

\begin{baihoc}
Thế giới không bao giờ thừa người tốt, nhất là người tốt mà còn có năng lực.
\textit{(The world never has too many good people, especially good people who are also capable.)}
\end{baihoc}

\begin{ghinhoanh}
\textbf{Short English to remember:}

Skill matters, but goodness is still necessary.

\end{ghinhoanh}

\begin{tuhocnhanh}
1. Vì sao người tốt luôn cần thiết? \textit{Why are good people always needed?}

2. Người tốt khác với người yếu ở đâu? \textit{How is goodness different from weakness?}

3. Em muốn dùng năng lực của mình theo kiểu nào? \textit{How do you want to use your abilities?}
\end{tuhocnhanh}
\ngancach

\section{Kiên Nhẫn Giúp Mình Đi Xa}
\begin{truyen}{Kiên Nhẫn Giúp Mình Đi Xa}{Patience Helps Us Go Far}
\chuhoa{K}{iên nhẫn không tạo kết quả ngay, nhưng nó ngăn em bỏ cuộc quá sớm.}
\textit{(Patience does not create instant results, but it keeps you from quitting too soon.)}

Trong đời sống, nhiều điều đáng giá đến chậm.
\textit{(Many valuable things in life arrive slowly.)}

Tình bạn tốt cần thời gian. Năng lực thật cần thời gian. Sự trưởng thành cũng cần thời gian.
\textit{(Good friendship, real ability, and maturity all require time.)}

Người thiếu kiên nhẫn hay đốt hết sức mình vào việc đòi kết quả ngay.
\textit{(Impatient people often burn themselves demanding immediate results.)}
\end{truyen}

\begin{baihoc}
Kiên nhẫn là cách con người cho điều tốt đủ thời gian để lớn lên.
\textit{(Patience is how we give good things enough time to grow.)}
\end{baihoc}

\begin{ghinhoanh}
\textbf{Short English to remember:}

Patience gives good things time to grow.

\end{ghinhoanh}

\begin{tuhocnhanh}
1. Vì sao kiên nhẫn quan trọng trong cuộc sống, không chỉ trong học tập? \textit{Why does patience matter in life, not only in study?}

2. Điều tốt nào cần thời gian để lớn? \textit{What good things require time to grow?}

3. Em đang thiếu kiên nhẫn với điều gì? \textit{What are you being impatient with?}
\end{tuhocnhanh}
\ngancach

\section{Sự Tử Tế Rất Mạnh}
\begin{truyen}{Sự Tử Tế Rất Mạnh}{Kindness Is Powerful}
\chuhoa{T}{ử tế thường bị nhầm là mềm yếu, nhưng thật ra nó rất mạnh.}
\textit{(Kindness is often mistaken for softness, but it is powerful.)}

Một lời dịu đúng lúc có thể giữ người khác lại với việc học, với tình bạn, thậm chí với chính họ.
\textit{(A gentle word at the right time can keep someone connected to learning, to friendship, even to themselves.)}

Tử tế không làm hết mọi vấn đề, nhưng nó làm bầu không khí con người sống với nhau bớt khắc nghiệt đi rất nhiều.
\textit{(Kindness does not solve everything, but it makes the atmosphere between people far less harsh.)}

Trong một thế giới dễ gắt, người giữ được sự tử tế là người rất có sức mạnh bên trong.
\textit{(In a harsh world, those who keep kindness carry inner strength.)}
\end{truyen}

\begin{baihoc}
Tử tế không phải thứ yếu. Nó là một sức mạnh giúp con người không làm nhau gãy thêm.
\textit{(Kindness is not secondary. It is a strength that keeps people from breaking one another further.)}
\end{baihoc}

\begin{ghinhoanh}
\textbf{Short English to remember:}

Kindness is a form of strength.

\end{ghinhoanh}

\begin{tuhocnhanh}
1. Vì sao sự tử tế mạnh hơn người ta tưởng? \textit{Why is kindness stronger than people think?}

2. Một lời tử tế đúng lúc có thể làm gì? \textit{What can a timely kind word do?}

3. Em có thể tử tế hơn với ai ngay hôm nay? \textit{With whom can you be kinder today?}
\end{tuhocnhanh}
\ngancach

\section{Học Để Hiểu Cuộc Sống}
\begin{truyen}{Học Để Hiểu Cuộc Sống}{Learn in Order to Understand Life}
\chuhoa{H}{ọc không chỉ để kiếm sống. Học còn để hiểu đời, hiểu người và hiểu chính mình.}
\textit{(We do not learn only to make a living. We also learn to understand life, other people, and ourselves.)}

Nếu chỉ giỏi việc mà không hiểu đời, con người vẫn có thể sống rất vụng.
\textit{(If we are only competent at work but do not understand life, we may still live clumsily.)}

Điều đẹp nhất của việc học là nó vừa cho em nghề, vừa cho em mắt nhìn.
\textit{(The most beautiful thing about learning is that it can give you both a profession and a way of seeing.)}

Khi đó, em không chỉ sống được. Em còn sống có chiều sâu hơn.
\textit{(Then you do not only survive. You live with more depth.)}
\end{truyen}

\begin{baihoc}
Việc học đẹp nhất là việc học giúp em vừa làm được việc, vừa hiểu đời cho ra đời.
\textit{(The best learning helps you both work well and understand life more truly.)}
\end{baihoc}

\begin{ghinhoanh}
\textbf{Short English to remember:}

Learn not only to work, but also to understand life.

\end{ghinhoanh}

\begin{tuhocnhanh}
1. Vì sao học không nên chỉ để kiếm sống? \textit{Why should learning not be only for earning a living?}

2. Hiểu đời giúp con người sống khác đi ra sao? \textit{How does understanding life change the way we live?}

3. Việc học của em đang mở thêm điều gì trong cách nhìn cuộc sống? \textit{What is your learning opening in the way you see life?}
\end{tuhocnhanh}
\ngancach
